\section{Experimental Setup and Supplementary Results}\label{app:exp}
We first describe the validation goals and experimental procedures, then present the setup and supplementary results in the order of Table~\ref{tab:exp-grid}.

\subsection{Validation Goals and Experiment Overview}\label{app:exp-overview}
The experiments in Section~\ref{sec:exp} validate the following conclusions. The statistical experiments (\exptag{STAT}) measure the lower-tail shape of the post-failure statistic (Theorems~\ref{thm:cross} and~\ref{thm:clt}) and the conditional miss rate and extra capacity of the two rules at finite size (Theorem~\ref{thm:coverage} and Corollary~\ref{cor:gauss}). For some configurations, we sample from the planted distribution in place of the failure-conditioned distribution. The planted sampling control (\exptag{PLANT}) tests the finite-size error of this substitution (Lemma~\ref{lem:plant}). The decoding-equivalence test (\exptag{DEC}) verifies on the full IBLT that random-sign encoding preserves decoding (Proposition~\ref{prop:decode}). The conditional-unbiasedness test (\exptag{UNB}) verifies that the post-failure estimate is conditionally unbiased given the mapping for every sign composition (Propositions~\ref{prop:moments} and~\ref{prop:composition}). The worst-configuration test (\exptag{WORST}) uses independent samples to test the conditional miss rate of the Skellam rule on the configurations closest to the risk upper bound, supplementing the statistical experiments. The protocol experiments (\exptag{PROTO}) measure the reliability and communication cost of the two-round protocol (Theorem~\ref{thm:protocol}). The computation-cost measurement (\exptag{COST}) measures the time and memory required for sizing and decoding.

\begin{table}[!htbp]
\centering
\caption{Experimental setup. Tags identify experiments in the main text and appendix. The last column lists the subsections reporting the results.}
\label{tab:exp-grid}
\small
\setlength{\tabcolsep}{3pt}
\begin{tabular}{@{}>{\raggedright\arraybackslash}p{.11\linewidth}>{\raggedright\arraybackslash}p{.17\linewidth}>{\raggedright\arraybackslash}p{.30\linewidth}>{\raggedright\arraybackslash}p{.24\linewidth}>{\raggedright\arraybackslash}p{.11\linewidth}@{}}
\toprule
Experiment & Goal & Parameter range & Samples and stopping rule & Location\\
\midrule
\exptag{STAT} statistical experiments & Lower-tail shape (Theorems~\ref{thm:cross}, \ref{thm:clt}); conditional miss rate and extra capacity (Theorem~\ref{thm:coverage}) & $k=3,4,5$, $M\in\{256,1024,4096,16384\}$, $c\in\{0.5,\dots,64\}$, load $\in\{0.95,\dots,3\}$ and $d=2,3,6$, load at most 3; $M=64$ boundary under the same rule & 189 + 42 configurations, about 20{,}000 failures each; fixed number of trials & \ref{sec:exp-setup}, \ref{sec:exp-dist}, \ref{sec:exp-cover}; \ref{app:sampling-impl}, \ref{app:limit}, \ref{app:cover-more}\\
\exptag{PLANT} planted sampling control & Finite-size error of the planted distribution (Lemma~\ref{lem:plant}) & Critical-load group: $k=3,4,5$, $M=64$--1024, load 0.85--0.95, direct failure probability at least $10^{-3}$; low-load group: for $k=3$, $M=1024$--16384, load 0.04--0.61; for $k=4,5$, $M=256,1024$, load 0.04--0.71 & Critical-load group: 21 configurations, 65{,}000 failures each for direct and planted; low-load group: 12 configurations, direct trial count fixed by $\lambda$, about 65{,}000 failures for $k=3$ and about 20{,}000 to 65{,}000 for $k=4,5$, 65{,}000 planted each & \ref{sec:exp-setup}; \ref{app:planted}\\
\exptag{DEC} decoding-equivalence test & Decoding equivalence (Proposition~\ref{prop:decode}) & Grid: $k\in\{3,4,5,6,10\}$, $M\in\{64,1024,16384\}$, load 0.5--3, five sign compositions; near threshold: $M\in\{1024,16384,65536\}$, load 0.90--1.05; rejected candidates: $M=16384$, load 0.5--1, nine sign compositions & Grid: 375 configurations, 500,000,000 paired trials in total; near threshold: 13,000,000; rejected candidates: 6,750,000 & \ref{sec:exp-decode}; \ref{app:decode-exp}\\
\exptag{UNB} conditional-unbiasedness test & Conditional unbiasedness (Propositions~\ref{prop:moments}, \ref{prop:composition}) & $k=3$--6, $M=64$--262144, critical window and its left side with load 0.4--0.95; Figures~\ref{fig:encoding}, \ref{fig:composition}: $M=16384$, failure probability about 0.1 & 203 configurations; Figure~\ref{fig:encoding}: 4,000,000 mappings, Figure~\ref{fig:composition}: 4 configurations & \ref{sec:exp-decode}; \ref{app:dist}\\
\exptag{WORST} worst-configuration test & Risk guarantee on the worst configurations (Theorem~\ref{thm:coverage}) & 12 configurations (Table~\ref{tab:confirm}); another 10 configurations for Figure~\ref{fig:scale} under the same design & 56{,}000--330{,}000 failures per configuration; stop when the block-ordered cumulative failure count reaches the target & \ref{sec:exp-cover}, \ref{sec:exp-e2e}; \ref{app:cover-more}\\
\exptag{PROTO} protocol experiments & Protocol reliability and communication cost (Theorem~\ref{thm:protocol}) & $k=3,4,5$, $M_1=256$--16384, load $\in\{0.95,1,1.05,1.5,3\}$, one-sided and balanced, $u\in\{0.01,0.05\}$ & 120 configurations, 13,940,000 trials in total per level of $u$ & \ref{sec:exp-e2e}; \ref{app:protocol}\\
\exptag{COST} computation-cost measurement & Time and memory of sizing and decoding & $M=16384$, $k=3,4,5$, $c$ from 4 to 314; protocol load 1.0--3 & 10 repetitions per item; candidate-check counts run separately, 30 trials per configuration & \ref{sec:exp-cost}; \ref{app:timing}\\
\bottomrule
\end{tabular}
\end{table}

Before launch, we fix the configurations, seeds, sampling sources, stopping rules, and budgets for the statistical experiments, worst-configuration test, planted sampling control, and protocol experiments using a single feasibility pilot. We fix the selection rule for the worst-configuration test before the statistical experiments begin, then derive the specific configurations and sample sizes from that rule after the statistical experiments finish (Appendix~\ref{app:cover-more}). We conduct the conditional-unbiasedness test and decoding-equivalence test in separate independent runs, and measure computation cost separately from the other runs (Appendix~\ref{app:timing}).

\subsection{Sampling Method}\label{app:sampling-impl}
For the statistical experiments, we choose the sampling source based on the cost estimated in the pilot. We use direct sampling when the estimated cost of obtaining 20{,}000 failures is at most 20 CPU core-hours, setting the number of trials to $\lceil2\times10^4/\widehat{\Pr}(F)\rceil$. For the remaining configurations, we use 20{,}000 planted samples. Figure~\ref{fig:exp-grid} shows each configuration's location in $c$ and load and its sampling source. We run the six points at load 3 ($d=\mathrm{round}(3\alpha_cM)$, with rounding making the load slightly above 3) separately with independent seeds after the statistical experiments finish. These points use the same sampling method and target failure count as the other configurations.

\begin{figure}[!htbp]
\centering\includegraphics[width=\linewidth]{../figure/fig-exp-grid.pdf}
\caption{The 189 configurations of the statistical experiments. Point shapes indicate the sampling source. Vertical lines mark the grouping boundaries of $c$, and horizontal shading marks load 0.95--1.05.}
\label{fig:exp-grid}
\end{figure}

Planted sampling selects a pair of difference keys uniformly at random and assigns them the same uniformly random $k$-cell set. We generate the remaining keys' mappings independently and uniformly, and draw the signs independently. This construction always produces a peeling failure. Direct sampling requires an average of $1/\Pr(F)$ mappings per failure, whereas planted sampling generates a mapping drawn from $Q$ directly. Lemma~\ref{lem:plant} gives its asymptotic relation to the true failure-conditioned distribution, and Appendix~\ref{app:planted} gives the finite-size comparison. Under iid random signs, the statistic is label-invariant, so we equivalently fix the last two labels as the duplicate edge pair in these experiments.

For the conditional-unbiasedness test, we combine direct sampling and splitting on an NVIDIA H100 with CPU sampling. The GPU direct sampler processes one graph per thread block and performs peeling in synchronized rounds. Each round first identifies degree-one cells, then uses atomic operations to claim and delete the corresponding edges, ensuring that each edge is deleted once. We use Philox to generate random numbers, organizing them by seed and trial index so that the random stream is independent of the GPU block partition. An independent CPU implementation re-peels the same exported batch of mappings and compares the final residual core graph by graph. For the initial direct sampler, we find no differences on 14,100 graphs, and the failure-probability statistical controls on all 15 configurations pass the preset criterion. Synchronized peeling and sequential CPU peeling yield the same final core. The CPU implementation handles path statistics that depend on the peeling order.

\subsection{Finite-Size Error of the Planted Distribution}\label{app:planted}
The planted sampling control (\exptag{PLANT}) compares finite-size failure samples from the planted distribution and the failure-conditioned distribution (Lemma~\ref{lem:plant}). We compare three quantities: the fraction of pure duplicate edge pairs among direct failures, the full-range KS distance between the two samples' post-failure $X$ distributions, and the miss/u of the Skellam rule on each sample. We compare the KS distance with the asymptotic 95\% critical value: 0.0075 with 65{,}000 failures on each side, and 0.011 with 20{,}000 on the direct side. This comparison is conservative because $X$ is discrete. We divide the control into two groups. The low-load group contains 12 configurations from the planted configurations of the statistical experiments and directly tests the planted samples used in the main text. The critical-load group contains 21 configurations at load 0.85 to 0.95, where direct failures are common, and tests how the error of the planted approximation varies with $M$. Figure~\ref{fig:planted} summarizes both groups.

Most of the 105 planted configurations in the statistical experiments lie far from the critical load. The expected number of duplicate edge pairs $\lambda=\binom d2/\binom Mk$ is therefore small, and Lemmas~\ref{lem:plant} and~\ref{lem:propL} imply that the failure probability is of the same order as $\lambda$. The low-load group consists of the 12 configurations for which direct sampling is affordable. Before sampling, we fix the direct-sampling trial count using $\lambda$ (Table~\ref{tab:v-low}). The four configurations with $k=3$ cover $c=1$ to 64 and load 0.04 to 0.61, with about 65{,}000 failures on each side. At $u=0.05$, all equivalence tests with bound $\pm0.1$ and one-sided level $0.05/4$ pass. At $u=0.01$, all difference intervals contain zero. The KS distances range from 0.0048 to 0.0079, with three below the critical value. The distance for $(3,4096,1024)$ is 0.0079 ($p=0.034$), so no difference is detectable after correction for 12 tests. The eight configurations with $k=4,5$ cover $c=0.5$ to 16 and load 0.04 to 0.71, where failures are rarer. Each has about 20{,}000 direct failures, except $(k,M,d)=(4,256,16)$ and $(5,256,128)$, which have about 65{,}000. Each has 65{,}000 planted samples. At $u=0.05$, all point estimates of the differences lie within $\pm0.04$, and all 95\% intervals contain zero. The equivalence tests with bound $\pm0.1$ and one-sided level $0.05/8$ pass for 5 configurations. For the other 3 configurations ($(4,1024,512)$, $(5,256,32)$, $(5,256,64)$), the interval endpoints lie within $\pm0.13$. At $u=0.01$, the difference intervals for 7 configurations contain zero. The difference for $(4,1024,512)$ is $-0.16$ (interval $[-0.30,-0.02]$), indicating higher miss in the planted sample. The KS distances range from 0.0048 to 0.0071, all below their respective critical values. Across the 12 configurations, the fraction of pure duplicate edge pairs among direct failures ranges from 0.988 to 1.000, consistent with the mechanism in Lemma~\ref{lem:plant}.

\begin{table}[!htbp]
\centering
\caption{Direct and planted sampling in the low-load group (Skellam rule, $u=0.05$). miss/u lists direct / planted. Difference gives the 95\% interval for direct minus planted. KS is the full-range KS distance of $X$, with 95\% critical value 0.0075 for 65{,}000 direct failures and 0.011 for 20{,}000.}
\label{tab:v-low}
\small
\begin{tabular}{@{}rrrrrrcccc@{}}
\toprule
& & & & & direct & pure duplicate edge pairs & & & \\
$k$ & $M$ & $d$ & $c$ & load & failure count & fraction & miss/u & difference & KS\\
\midrule
3 & 1024 & 32 & 1 & 0.04 & 65061 & 1.000 & 0.94 / 0.94 & $[-0.05,+0.04]$ & 0.0059\\
3 & 4096 & 512 & 8 & 0.15 & 64977 & 1.000 & 0.95 / 0.96 & $[-0.06,+0.03]$ & 0.0048\\
3 & 4096 & 1024 & 16 & 0.31 & 65136 & 0.999 & 0.92 / 0.94 & $[-0.07,+0.03]$ & 0.0079\\
3 & 16384 & 8192 & 64 & 0.61 & 65530 & 0.999 & 0.98 / 0.99 & $[-0.06,+0.04]$ & 0.0051\\
\addlinespace
4 & 256 & 8 & 0.5 & 0.04 & 20071 & 0.998 & 0.54 / 0.52 & $[-0.03,+0.07]$ & 0.0071\\
4 & 256 & 16 & 1 & 0.08 & 65839 & 0.995 & 0.68 / 0.69 & $[-0.05,+0.03]$ & 0.0057\\
4 & 1024 & 128 & 4 & 0.16 & 20067 & 0.997 & 0.95 / 0.95 & $[-0.06,+0.07]$ & 0.0049\\
4 & 1024 & 256 & 8 & 0.32 & 20265 & 0.994 & 0.98 / 0.97 & $[-0.06,+0.08]$ & 0.0048\\
4 & 1024 & 512 & 16 & 0.65 & 20038 & 0.988 & 0.94 / 0.97 & $[-0.10,+0.04]$ & 0.0062\\
5 & 256 & 32 & 2 & 0.18 & 20137 & 0.999 & 1.00 / 0.98 & $[-0.05,+0.09]$ & 0.0068\\
5 & 256 & 64 & 4 & 0.36 & 19885 & 0.997 & 0.91 / 0.95 & $[-0.11,+0.03]$ & 0.0048\\
5 & 256 & 128 & 8 & 0.71 & 65470 & 0.995 & 0.99 / 1.01 & $[-0.07,+0.02]$ & 0.0063\\
\bottomrule
\end{tabular}
\end{table}

Direct sampling is infeasible for the remaining planted configurations. Among the 77 planted configurations with $k=4,5$, 76 have $\lambda$ below $10^{-6}$, with a minimum of about $10^{-19}$. The feasibility pilot yields no failures in 68 of 75 configurations. Collecting 65{,}000 failures requires at least $6.5\times10^{10}$ trials, and about half the configurations require more than $10^{15}$, exceeding the feasible range of direct sampling on CPU and GPU. For these configurations, Lemma~\ref{lem:plant} controls the distance between the planted distribution and the failure-conditioned distribution, with an error term that tends to zero when $d=o(M)$. The approximation is most accurate precisely where direct sampling is hardest.

The critical-load group tests the range of applicability of the planted approximation. The statistical experiments include five planted configurations at load 0.95, all with $M\ge4096$. This group uses smaller $M$ to compare the two samples where larger stopping sets cause most failures. We determine the sample size using an equivalence test on the miss/u difference at $u=0.05$, with bound $\pm0.1$, one-sided level $0.05/21$, and 80\% power when the difference is zero. This gives 65{,}000 failures on each side. Table~\ref{tab:v} reports the results for these 21 configurations.

\begin{table}[!htbp]
\centering
\caption{Direct and planted sampling in the critical-load group (Skellam rule). miss/u lists direct / planted. KS is the full-range KS distance of $X$, with 95\% critical value 0.0075.}
\label{tab:v}
\small
\begin{tabular}{@{}rrrrcccc r@{}}
\toprule
$k$ & $M$ & load & pure duplicate edge pairs & \multicolumn{2}{c}{$u=0.01$} & \multicolumn{2}{c}{$u=0.05$} & KS\\
\cmidrule(lr){5-6}\cmidrule(lr){7-8}
& & & fraction & miss/u & difference & miss/u & difference & \\
\midrule
3 & 64 & 0.86 & 0.09 & 0.44 / 0.48 & $-0.04\pm0.07$ & 0.87 / 0.88 & $-0.01\pm0.04$ & 0.0026\\
3 & 64 & 0.90 & 0.05 & 0.49 / 0.52 & $-0.03\pm0.08$ & 0.82 / 0.87 & $-0.05\pm0.04$ & 0.0085\\
3 & 64 & 0.95 & 0.02 & 0.44 / 0.51 & $-0.07\pm0.07$ & 0.87 / 0.94 & $-0.07\pm0.05$ & 0.0088\\
4 & 64 & 0.85 & 0.01 & 0.56 / 0.49 & $+0.08\pm0.08$ & 0.89 / 0.92 & $-0.03\pm0.05$ & 0.0105\\
4 & 64 & 0.89 & 0.01 & 0.54 / 0.65 & $-0.11\pm0.08$ & 0.89 / 0.96 & $-0.07\pm0.05$ & 0.0097\\
4 & 64 & 0.95 & 0.00 & 0.39 / 0.55 & $-0.16\pm0.07$ & 0.81 / 0.93 & $-0.13\pm0.04$ & 0.0131\\
5 & 64 & 0.85 & 0.00 & 0.58 / 0.69 & $-0.11\pm0.09$ & 0.96 / 1.07 & $-0.11\pm0.05$ & 0.0130\\
5 & 64 & 0.89 & 0.00 & 0.57 / 0.69 & $-0.12\pm0.09$ & 1.00 / 1.13 & $-0.13\pm0.05$ & 0.0142\\
5 & 64 & 0.96 & 0.00 & 0.55 / 0.72 & $-0.16\pm0.09$ & 0.89 / 1.04 & $-0.15\pm0.05$ & 0.0118\\
3 & 256 & 0.85 & 0.63 & 0.75 / 0.62 & $+0.13\pm0.09$ & 0.89 / 0.88 & $+0.01\pm0.04$ & 0.0063\\
3 & 256 & 0.90 & 0.13 & 0.75 / 0.66 & $+0.09\pm0.09$ & 0.95 / 0.88 & $+0.07\pm0.05$ & 0.0080\\
3 & 256 & 0.95 & 0.02 & 0.77 / 0.64 & $+0.14\pm0.09$ & 0.92 / 0.91 & $+0.01\pm0.05$ & 0.0060\\
4 & 256 & 0.90 & 0.01 & 0.71 / 0.70 & $+0.01\pm0.09$ & 0.93 / 0.98 & $-0.05\pm0.05$ & 0.0045\\
4 & 256 & 0.95 & 0.00 & 0.72 / 0.73 & $-0.01\pm0.09$ & 0.93 / 0.97 & $-0.05\pm0.05$ & 0.0052\\
5 & 256 & 0.90 & 0.00 & 0.71 / 0.75 & $-0.04\pm0.09$ & 0.92 / 0.94 & $-0.02\pm0.05$ & 0.0052\\
5 & 256 & 0.95 & 0.00 & 0.73 / 0.80 & $-0.07\pm0.09$ & 0.95 / 0.94 & $+0.00\pm0.05$ & 0.0027\\
3 & 1024 & 0.85 & 0.97 & 0.79 / 0.78 & $+0.01\pm0.10$ & 0.92 / 0.91 & $+0.02\pm0.05$ & 0.0033\\
3 & 1024 & 0.90 & 0.96 & 0.83 / 0.86 & $-0.03\pm0.10$ & 0.93 / 0.95 & $-0.02\pm0.05$ & 0.0037\\
3 & 1024 & 0.95 & 0.08 & 0.84 / 0.84 & $+0.01\pm0.10$ & 0.96 / 0.96 & $-0.00\pm0.05$ & 0.0091\\
4 & 1024 & 0.95 & 0.00 & 0.89 / 0.91 & $-0.02\pm0.10$ & 0.97 / 0.94 & $+0.03\pm0.05$ & 0.0054\\
5 & 1024 & 0.95 & 0.00 & 0.85 / 0.79 & $+0.06\pm0.10$ & 0.96 / 0.95 & $+0.01\pm0.05$ & 0.0039\\
\bottomrule
\end{tabular}
\end{table}

For all five tested $M=1024$ configurations, the equivalence tests with bound $\pm0.1$ pass at $u=0.05$, and the difference intervals at both levels of $u$ contain zero. These include configurations where the fraction of pure duplicate edge pairs among direct failures is 0.97 and others where it is close to zero. For seven of the nine $M=64$ configurations, the tests do not establish that the difference lies within the bound. The planted sample has higher miss in all seven, with the difference interval extending as far as $-0.20$. The tests also do not establish the bound for three of the seven $M=256$ configurations: at $k=3$ and load 0.90 the direct sample has higher miss, with difference $+0.07$; at $k=4$ and loads 0.90 and 0.95 the difference is $-0.05$ in both. For the two configurations with $k=3$, $M=256$, and loads 0.85 and 0.95, both difference intervals contain zero at $u=0.05$. At $u=0.01$, the direct sample's miss/u is higher by 0.13 (interval $[+0.04,+0.22]$) and 0.14 (interval $[+0.04,+0.23]$), respectively.

The KS distance decreases as $M$ grows. For the five $M=1024$ configurations, it ranges from 0.0033 to 0.0091, with four below the critical value. The configuration $(k,M,d)=(3,1024,796)$ exceeds the critical value because of differences in the middle of the distribution. The maximum difference in the lower tail between the 0.2\% and 15\% quantiles of the pooled sample is only 0.0039. For the seven $M=256$ configurations, the distance ranges from 0.0027 to 0.0080, with six below the critical value. For the nine $M=64$ configurations, it ranges from 0.0026 to 0.0142, with eight above the critical value. At $M=1024$, the configuration with a pure duplicate edge pair fraction of 0.97 and those with fractions close to zero have similarly small KS distances. Thus, the distributions of $X$ can be close even when the failure mechanisms differ. The fraction of pure duplicate edge pairs alone is insufficient to predict the distribution error of $X$, which also depends on overlap fluctuations on the $s(d)$ scale.

\begin{figure}[!htbp]
\centering\includegraphics[width=\linewidth]{../figure/fig-planted.pdf}
\caption{Direct versus planted sampling. The upper three panels show the critical-load group, grouped by $M$; the bottom panel shows the low-load group. The vertical axis labels $k$ and load. Left: Skellam rule miss/u difference (direct minus planted) with 95\% intervals and shading at $\pm0.1$. Right: full-range KS distance of $X$; gray ticks mark the 95\% critical value for each sample size.}
\label{fig:planted}
\end{figure}

\subsection{Random-Sign Encoding Does Not Change Decoding}\label{app:decode-exp}
The decoding-equivalence test (\exptag{DEC}) uses the two-round protocol's actual IBLT encoder and decoder. We build and subtract the two parties' tables, then compare first-round decoding under standard encoding and random-sign encoding. The test consists of three groups of paired trials, all using the same construction and verification procedure.

In each trial, we generate distinct keys and fresh seeds for the cell hash, sign hash, and checksum hash. Both encodings share the same keys and hashes. Before decoding, we verify cell by cell that the three fields after cancellation of the common elements match those obtained by encoding the difference directly. An independent reference uses only the unsigned incidence graph of the true difference and computes its residual 2-core by the degree-one rule. Whenever the actual decoder accepts a key, we check the true residual degree, key, and direction, then update the reference state. These checks cover trajectories ending in both success and failure. After termination, we verify that the recovered keys are exactly the difference outside the 2-core, with all directions correct. We reconstruct the three fields from the true residual, compare them cell by cell, and check that all fields are zero if and only if the residual core is empty.

The first group is a grid spanning the full load range: $k\in\{3,4,5,6,10\}$, $M\in\{64,1024,16384\}$, load $d/(\alpha_cM)\in\{0.5,0.95,1,1.05,3\}$, and positive-direction ratio $d_+/d$ with target values $0,1/4,1/2,3/4,1$. We round $d$ and $d_+$ to integers from the target values. We use the corresponding peeling threshold for each $k$, giving 375 configurations in total. The number of common elements is $\min\{1000,M\}$. For $M=64,1024,16384$, we run 3,000,000, 900,000, and 100,000 trials per configuration, respectively. This gives 500,000,000 trials in total, with 100,000,000 for each $k$.

Table~\ref{tab:decode-paired} reports the results by $k$. Across all configurations, we observe zero paired success/failure disagreements, non-pure acceptances, incorrect recoveries, or residual-core disagreements. The two decoders together check 211,516,840,988 acceptance operations, all satisfying the true degree-one and direction conditions.

The second group samples loads more densely near the threshold and provides the results in Table~\ref{tab:decode} of the main text. The configurations use $k\in\{3,4,5,6,10\}$, $M\in\{1024,16384,65536\}$, 10 load values between 0.90 and 1.05, and $d_-/d\in\{0,1/2\}$. For $M=1024,16384,65536$, we run 100,000, 20,000, and 10,000 trials per configuration, respectively, for a total of 13,000,000. All checks pass. At $M=16384$, the failure probability ranges from 0.020 to 0.060 at load 0.99 and from 0.993 to 0.998 at load 1.01. At $M=65536$, almost no trials fail at load 0.99, and almost all fail at load 1.01. After failure, the decoder still recovers the difference outside the 2-core. At $M=65536$ and load 1.02, the average recovery fraction ranges from 0.18 ($k=10$) to 0.50 ($k=3$).

The third group of paired trials counts rejected candidates separately for each encoding. These are cells with count $\pm1$ that fail the checksum or cell-ownership check. The configurations use $M=16384$, $k\in\{3,4,5,6,10\}$, load $0.5,0.95,1$, and nine values of $d_-/d$ from 0 to 1 in steps of $1/8$. We run 50,000 trials per configuration, for a total of 6,750,000, and all checks pass. Figure~\ref{fig:decode} reports the results for $k=3$ to 6. Standard encoding has no rejected candidates under one-sided differences because the count equals the degree. As the difference approaches balanced, more opposite-sign cancellation occurs and more candidates are rejected. Under random-sign encoding, the number of rejected candidates differs by at most 0.3\% across the nine sign compositions and matches standard encoding in the balanced case. This agrees with Proposition~\ref{prop:composition}: the coding coefficient is an iid random sign independent of $(d_+,d_-)$. The results for $k=10$ follow the same pattern, with both encodings yielding about 1.27 rejected candidates per difference key at load 1. Across all three groups, the two decoders together reject $2.8\times10^{11}$ non-pure candidates with count $\pm1$, with zero false acceptances. The cell-ownership check rejects zero candidates, and the checksum rejects every incorrect candidate.

\begin{figure}[!htbp]
\centering\includegraphics[width=\linewidth]{../figure/fig-decode.pdf}
\caption{Rejected candidates per first-round decoding, divided by $d$ ($M=16384$; both encodings share the same mappings for each configuration). Solid lines with filled markers show standard encoding; dashed lines with open markers show random-sign encoding.}
\label{fig:decode}
\end{figure}

\begin{table}[!htbp]
\centering
\caption{First-round paired decoding results.}
\label{tab:decode-paired}
\begin{tabular}{rrrr}
\toprule
$k$ & paired trials & successes & failures\\
\midrule
3 & 100,000,000 & 36,532,515 & 63,467,485\\
4 & 100,000,000 & 35,900,798 & 64,099,202\\
5 & 100,000,000 & 33,530,250 & 66,469,750\\
6 & 100,000,000 & 32,553,378 & 67,446,622\\
10 & 100,000,000 & 30,233,256 & 69,766,744\\
\bottomrule
\end{tabular}
\end{table}

\subsection{The Post-Failure Estimate Is Conditionally Unbiased for Every Sign Composition}\label{app:dist}
Of the 203 configurations in the conditional-unbiasedness test (\exptag{UNB}), 196 have at least 1000 failure samples. Under random-sign encoding, $|\E[X\mid F]|$ is at most 0.032 for these configurations, with 93\% within two standard errors of zero, consistent with the exact unbiasedness in Proposition~\ref{prop:moments}. The conditional variance to the left of the critical window is 0.972--1.049. Among the 84 direct-sampling configurations in the statistical experiments, $|\E[X\mid F]|$ under random-sign encoding is at most 0.019, with 98\% within two standard errors of zero. In the statistical experiments at load 1.0 and failure probability about 0.67, standard encoding with one-sided differences has conditional means of $+0.20$, $+0.02$, and $-0.08$ for $k=3,4,5$, respectively (standard error 0.007).

Table~\ref{tab:mean-by-k} reports the values in Figure~\ref{fig:composition} of the main text. For each $(k,M,d)$, we record five sign compositions under both encodings. The ten datasets share the same mappings, and both encodings use all failure records. The four configurations have failure probabilities of 0.099--0.102. All 20 values under random-sign encoding have absolute value at most 0.005 and lie within two standard errors of zero. We generate the data for Figure~\ref{fig:encoding} using the sampling code from the interactive experiment page Peeling Lens. The code computes $X$ for three encodings on each random mapping, drawing random signs from an independent generator. This configuration has $c\approx104$. Under standard encoding, $\E[X\mid F]=+0.599\pm0.001$ for one-sided differences and $-0.002$ for balanced differences. Under random-sign encoding, it is $+0.001\pm0.001$, with conditional variance 1.007. The failure samples and the full sample have similar distributions, consistent with the large-$c$ behavior in Theorem~\ref{thm:clt}.

\begin{table}[!htbp]
\centering
\caption{Values in Figure~\ref{fig:composition}: mean post-failure $X$ ($M=16384$, failure probability about 0.1). Standard errors are 0.0018 for $k=3$ and 0.0031--0.0032 for $k=4,5,6$.}
\label{tab:mean-by-k}
\begin{tabular}{rrlrrrrr}
\toprule
& & & \multicolumn{5}{c}{$d_-/d$}\\
\cmidrule(lr){4-8}
$k$ & $d$ & encoding & 0 (one-sided) & 0.1 & 0.25 & 0.4 & 0.5 (balanced)\\
\midrule
3 & 13306 & standard & $+0.638$ & $+0.405$ & $+0.158$ & $+0.028$ & $+0.003$\\
  &       & random sign & $+0.002$ & $-0.001$ & $0.000$ & $-0.004$ & $+0.002$\\
\addlinespace
4 & 12567 & standard & $+0.084$ & $+0.052$ & $+0.016$ & $-0.002$ & $-0.003$\\
  &       & random sign & $-0.001$ & $-0.002$ & $+0.005$ & $+0.001$ & $-0.001$\\
\addlinespace
5 & 11418 & standard & $-0.218$ & $-0.137$ & $-0.054$ & $-0.009$ & $+0.002$\\
  &       & random sign & $+0.003$ & $-0.005$ & $0.000$ & $-0.002$ & $-0.004$\\
\addlinespace
6 & 10363 & standard & $-0.387$ & $-0.244$ & $-0.096$ & $-0.021$ & $-0.004$\\
  &       & random sign & $+0.002$ & $+0.003$ & $+0.005$ & $+0.001$ & $+0.001$\\
\bottomrule
\end{tabular}
\end{table}

\subsection{Distance between the Post-Failure Lower Tail and the Limit Distribution}\label{app:limit}
We use the failure records from the statistical experiments (\exptag{STAT}). For Figure~\ref{fig:crosscdf}, we compute $\widetilde d=\sum_vC_v^2/k$ from the raw integer sums of squares, retain all empirical values, and standardize them on the same scale as $G_d$ in the main text. We compute the KS distance immediately to the left and right of every value taken by either distribution. By \eqref{eq:raw-decomp} in Appendix~\ref{app:ams}, the difference between $\widehat d$ and $\widetilde d$ is about $0.02s(d)$ at $M=16384$. This perturbation is small relative to the standard deviation, yet it shifts the discrete atoms. We therefore use $\widetilde d$, as in Theorem~\ref{thm:cross}, to validate the distribution.

Table~\ref{tab:limit-ks} reports the full KS distances computed from the raw integer records of the statistical experiments. The column for the discrete limit distribution compares $\widetilde d$ with $G_d$ for $c=1,2,4$ and $k=3,4,5$, giving nine configurations for each $M$. We use direct and planted samples for $M\le1024$ and planted samples for $M\ge4096$. The Gaussian control in the same column uses $(\widetilde d-d)/s(d)$. The Gaussian column on the right compares $X=(\widehat d-d)/s(d)$ with $N(0,1)$ for the 75 direct-sampling configurations with $c\ge8$, including the critical window and overload. Each configuration has about 20{,}000 failures. At this fixed sample size, the per-configuration 95\% DKW error bound is about 0.010. The discrete approximation has visible finite-size error for small tables, which generally decreases as $M$ grows. The Gaussian quantile error also decreases as $M$ grows.

\begin{table}[!htbp]
\centering
\caption{Full post-failure KS distance. The discrete limit uses $\widetilde d$; the Gaussian uses $\widehat d$. Quantile error is in units of $s(d)$ and gives the maximum absolute error across configurations.}
\label{tab:limit-ks}
\small
\begin{tabular}{rcccccc}
\toprule
& \multicolumn{3}{c}{discrete limit ($c=1,2,4$)} & \multicolumn{3}{c}{Gaussian ($c\ge8$, direct sampling)}\\
\cmidrule(lr){2-4}\cmidrule(lr){5-7}
$M$ & KS & KS of Gaussian ($c=1$) & configs & KS & 1\% quantile error & configs\\
\midrule
256 & 0.019--0.034 & 0.237--0.252 & 9 & 0.010--0.021 & 0.18 & 23\\
1024 & 0.008--0.023 & 0.224--0.244 & 9 & 0.005--0.015 & 0.10 & 23\\
4096 & 0.004--0.016 & 0.215--0.230 & 9 & 0.005--0.014 & 0.06 & 17\\
16384 & 0.003--0.012 & 0.214--0.230 & 9 & 0.003--0.010 & 0.05 & 12\\
\bottomrule
\end{tabular}
\end{table}

The conditional variance reflects the same structure. Among the 84 direct-sampling configurations in the statistical experiments, the conditional variance under random-sign encoding is 0.98--1.06 for $c\ge8$ and 0.985--1.020 within the critical window. For small $c$, it is well above one: about 1.12 at $c=4$, about 1.5 at $c=2$, and above 10 when $c<1$. These values agree with $1+2/c^2$, obtained by converting the limit distribution's variance $4+2c^2$ to the $X$ scale. The excess comes from the duplicate edge pair's $\pm2$ shift, which is non-negligible relative to $s(d)$.

\subsection{Conditional Miss Rate and Extra Capacity at Finite Size}\label{app:cover-more}
Of the 189 configurations in the statistical experiments, 84 use direct sampling and 105 use planted sampling. The Skellam rule's maximum miss/u on the direct-sampling configurations is 1.04 at both levels of $u$, occurring within the critical window. Its maximum on the planted configurations is 1.03, computed under the planted distribution. Each configuration has about 20{,}000 failures, corresponding to about 200 misses at $u=0.01$. The standard error of miss/u is about 0.07, so a value slightly above 1, such as 1.04, is statistically indistinguishable from 1. On the low-load configurations with $c<8$, the Skellam rule's median miss/u at $u=0.01$ is at most 0.76. Every sample yields at least one accepted candidate. If no candidate is accepted, $U=+\infty$, which does not count as a miss. For $c\ge8$, the Gaussian rule gives miss/u of 0.83--1.19 across all configurations. At $c=1$, $u=0.01$, it is about 11 for $k=3$ and about 8.5 for $k=4,5$. When $c$ is small and the observation is near $d-2$, observations for $d=2,3$ fall below every candidate's acceptance threshold. The Gaussian rule then returns $+\infty$, giving no miss. On the direct-sampling configurations with $c\ge8$ in the statistical experiments, the Skellam rule's median miss/u at the two levels of $u$ is 0.79--0.86 and 0.95--0.97.

At $u=0.01$ and $c\ge8$, the Skellam rule's median extra capacity is 25\%, 11.3\%, 5.4\%, and 2.6\% for $M=256$, 1024, 4096, 16384, respectively. The corresponding 90\% quantiles are at most 41\%, 18\%, 8.5\%, and 4.1\%. The Gaussian rule's medians are 23.5\%, 10.8\%, 5.3\%, and 2.6\%, differing from the Skellam rule by at most 1.5 percentage points. At $c=1,2,4$, the median of $U-d$ is 5, 8, 15 for the Skellam rule and 3, 6, 14 for the Gaussian rule. For fixed small $d$, the 90\% quantile of the upper bound at $d=2$ exceeds $d$ by 4--5. Among the 42 configurations with $M=64$, the Skellam rule's maximum miss/u is 1.08 at $u=0.05$ ($k=4$, $d=6$, 95\% Wilson interval 1.02--1.15) and 1.04 at $u=0.01$.

Table~\ref{tab:gauss-c} reports the Gaussian rule's miss/u at fixed $c$ for $M=16384$, corresponding to the triangles with $c\le16$ in Figure~\ref{fig:coverage} of the main text. The excess grows rapidly as $c$ decreases. The miss/u is about 8--11 at $c=1$ and above 33 at $c=0.5$ for $u=0.01$. For $c\le4$, the values at $M=4096$ and 16384 differ by at most 8\%, so the excess persists as the table size grows.

\begin{table}[!htbp]
\centering
\caption{Gaussian rule miss/u at fixed $c$ ($M=16384$).}
\label{tab:gauss-c}
\begin{tabular}{rrrrrrr}
\toprule
& \multicolumn{3}{c}{$u=0.01$} & \multicolumn{3}{c}{$u=0.05$}\\
\cmidrule(lr){2-4}\cmidrule(lr){5-7}
$c$ & $k=3$ & $k=4$ & $k=5$ & $k=3$ & $k=4$ & $k=5$\\
\midrule
0.5 & 36.4 & 33.2 & 33.2 & 9.37 & 8.81 & 9.39\\
1 & 11.3 & 8.74 & 8.27 & 4.09 & 4.24 & 4.60\\
2 & 3.06 & 2.48 & 2.60 & 2.00 & 1.70 & 1.94\\
4 & 1.51 & 1.41 & 1.50 & 1.21 & 1.19 & 1.20\\
8 & 1.12 & 1.03 & 1.08 & 1.07 & 1.08 & 1.06\\
16 & 0.98 & 0.93 & 0.99 & 0.98 & 1.01 & 1.00\\
\bottomrule
\end{tabular}
\end{table}


The worst-configuration test (\exptag{WORST}) uses samples independent of the statistical experiments. We fix the selection rule and sample sizes before sampling. We first select the five configurations whose miss/u exceeds 1 in exploration. We fill the remaining slots by ranking configurations in descending order of the larger one-sided 95\% exact upper bound divided by $u$ over the two levels of $u$ in the statistical experiments. Four highly ranked configurations have only planted samples, and direct sampling cannot obtain failures at feasible cost. Following the rule, we skip these in favor of the next configurations. Their miss/u under $Q$ is at most 1.02. For each configuration and each level of $u$, we use the previously estimated 95\% Wilson upper bound (at most 1) as the planning value. We choose the failure count to give a passing probability of at least $1-0.2/24$ for each judgment and take the larger count over the two levels. At the planning value, this gives at least 80\% probability that all 24 judgments pass. A judgment passes if its one-sided simultaneous confidence upper bound is at most $1.1u$. We use the Clopper--Pearson interval with a Bonferroni correction over the 24 judgments. All 24 pass, with a maximum upper bound of $1.024u$. Of the 24 two-sided 95\% intervals, 22 lie entirely below 1 and the other 2 cover 1. Only one test configuration uses $M=256$. We sample in blocks and analyze all blocks through the first point at which the cumulative failure count in block order reaches the target. Stopping depends only on the failure count. For the two test configurations with $k=3$, $M=1024$, the Gaussian rule's miss/u is 11.1 ($u=0.01$) and 4.1 ($u=0.05$) at $c=1$, and 2.9 and 1.8 at $c=2$. Table~\ref{tab:confirm} reports the load, failure count, miss/u, and one-sided simultaneous confidence upper bound for each test configuration. The segments in Figure~\ref{fig:confirm} of the main text are two-sided 95\% intervals.

\begin{table}[!htbp]
\centering
\caption{Skellam rule miss/u on the \exptag{WORST} configurations, with one-sided simultaneous confidence upper bounds in parentheses (per-item level $0.05/24$).}
\label{tab:confirm}
\begin{tabular}{lrrcc}
\toprule
configuration & load & failures & $u=0.01$ & $u=0.05$\\
\midrule
$k=3$, $M=1024$, $d=32$ & 0.04 & 55{,}509 & 0.53 (0.62) & 0.97 (1.02)\\
$k=3$, $M=1024$, $d=64$ & 0.08 & 289{,}001 & 0.90 (0.95) & 0.91 (0.94)\\
$k=3$, $M=1024$, $d=192$ & 0.23 & 289{,}394 & 0.79 (0.83) & 0.95 (0.97)\\
$k=4$, $M=256$, $d=128$ & 0.65 & 105{,}689 & 0.74 (0.82) & 0.93 (0.97)\\
$k=4$, $M=1024$, $d=751$ & 0.95 & 298{,}439 & 0.85 (0.90) & 0.95 (0.98)\\
$k=3$, $M=4096$, $d=3352$ & 1.00 & 333{,}476 & 0.93 (0.98) & 0.97 (1.00)\\
$k=5$, $M=4096$, $d=3018$ & 1.05 & 332{,}350 & 0.90 (0.95) & 0.99 (1.01)\\
$k=5$, $M=1024$, $d=755$ & 1.05 & 332{,}353 & 0.83 (0.87) & 0.96 (0.98)\\
$k=4$, $M=16384$, $d=16384$ & 1.29 & 293{,}250 & 0.95 (1.01) & 0.98 (1.01)\\
$k=3$, $M=4096$, $d=5029$ & 1.50 & 332{,}350 & 0.91 (0.96) & 0.98 (1.00)\\
$k=4$, $M=16384$, $d=18980$ & 1.50 & 290{,}808 & 0.96 (1.02) & 0.98 (1.00)\\
$k=4$, $M=16384$, $d=37959$ & 3.00 & 289{,}952 & 0.95 (1.01) & 0.97 (1.00)\\
\bottomrule
\end{tabular}
\end{table}

Of the 12 configurations in the upper row of Figure~\ref{fig:scale}, two ($k=3$, $M=4096$ and $k=4$, $M=16384$) are test configurations. The other 10 follow the same design, with sample size and stopping rule fixed before sampling, and each has about 290{,}000 direct-sampling failures. Under random-sign encoding, the common elements cancel in the difference table and sign composition does not affect the post-failure distribution. The $|b|\le1$ used by the protocol under balanced differences also leaves the upper bound unchanged. The protocol experiments, statistical experiments, and these direct samples therefore estimate the same conditional miss rate, and we pool them by configuration in the figure. For these 10 configurations, the Skellam rule's one-sided upper bound is at most $1.025u$ at $u=0.01$ and $1.003u$ at $u=0.05$.

\subsection{Reliability and Communication Cost of the Two-Round Protocol}\label{app:protocol}
In each trial, we generate the common elements and difference elements separately. For one-sided differences, all differences lie on the $A$ side; for balanced differences, half lie on each side. Common elements cancel in the difference table and participate as usual in table construction and cell-by-cell subtraction. Trials that succeed in the first round require no upper bound, so only first-round-failure trials enter the comparison. We evaluate the two levels of $u$ separately on the same batch of trials. We impose a minimum of 64 cells on the second-round table length specified in Section~\ref{sec:flow}. The second round decodes independently on the new table, consistent with Theorem~\ref{thm:protocol}. All second-round recovery results match the true difference element by element. The rounds use independent cell hashes and share the sign hash and checksum hash. We count 20 bytes per cell (4 bytes for the transmitted count and 8 bytes each for keySum and checksum), plus an 8-byte global counter for random-sign encoding. All upper bounds use the same keys and hashes. For each first-round-failure trial and each bound, we record $U_b$, the second-round size, and the second-round outcome, pairing comparisons by trial. Both levels of $u$ use the same number of trials per configuration, set to $\lceil2\times10^4/\widehat{\Pr}(F_1)\rceil$ with a cap of 3,000,000. The 13 configurations with no first-round failure in the pilot receive a fixed 100,000 trials. Among them, the eight configurations with $k=4,5$ at load 0.95 and $M_1\ge4096$ still have no first-round failure in these trials. These configurations contribute only to the failure-rate and byte statistics. We compute intervals for the success-rate difference using the Newcombe hybrid score method for paired proportion differences.

Table~\ref{tab:e2e} summarizes the protocol results from Section~\ref{sec:exp-e2e} of the main text by $M_1$ and $u$. At $M_1=256$, the final two-round failure rate is at most 0.3\%. The conditional miss rate on balanced configurations agrees with the statistical experiments in Section~\ref{sec:exp-cover}. Second-round failure after $U_b\ge d$ is 0.06--0.08\% at $M_1=256$. The oracle uses a smaller second-round table. In paired comparisons by trial, the Skellam rule's success rate is up to 0.35 percentage points higher than the oracle's on some $M_1=256$ configurations.

\begin{table}[!htbp]
\centering
\caption{Two-round protocol summary ($k=3,4,5$). miss/u gives the median / maximum over balanced configurations with at least 1000 first-round failures. Final failure is the maximum over configurations. Second-round size and total bytes are medians relative to the oracle. The oracle is independent of $u$ and appears only in the upper half.}
\label{tab:e2e}
\begin{tabular}{@{}llrrrr@{}}
\toprule
$M_1$, $u$ & upper bound & miss/u & final failure & second-round size & total bytes\\
\midrule
256, 0.01 & Skellam & 0.73 / 0.87 & 0.24\% & +25.6\% & +14.7\%\\
& Gaussian & 1.00 / 1.12 & 0.25\% & +23.9\% & +13.7\%\\
& oracle & -- & 0.51\% & 0 & 0\\
1024, 0.01 & Skellam & 0.88 / 1.00 & 0.06\% & +11.4\% & +6.5\%\\
& Gaussian & 1.04 / 1.17 & 0.08\% & +11.1\% & +6.3\%\\
& oracle & -- & 0.09\% & 0 & 0\\
4096, 0.01 & Skellam & 0.92 / 1.01 & 0.02\% & +5.4\% & +3.1\%\\
& Gaussian & 1.01 / 1.09 & 0.03\% & +5.3\% & +3.0\%\\
& oracle & -- & 0.03\% & 0 & 0\\
16384, 0.01 & Skellam & 0.96 / 1.07 & 0.01\% & +2.6\% & +1.5\%\\
& Gaussian & 1.02 / 1.13 & 0.01\% & +2.6\% & +1.5\%\\
& oracle & -- & 0.02\% & 0 & 0\\
\addlinespace
256, 0.05 & Skellam & 0.93 / 1.01 & 0.29\% & +16.9\% & +9.7\%\\
& Gaussian & 1.01 / 1.07 & 0.27\% & +16.3\% & +9.4\%\\
1024, 0.05 & Skellam & 0.95 / 1.05 & 0.13\% & +7.8\% & +4.5\%\\
& Gaussian & 1.02 / 1.09 & 0.11\% & +7.7\% & +4.4\%\\
4096, 0.05 & Skellam & 0.97 / 1.04 & 0.02\% & +3.8\% & +2.2\%\\
& Gaussian & 1.01 / 1.07 & 0.02\% & +3.8\% & +2.1\%\\
16384, 0.05 & Skellam & 1.00 / 1.05 & 0.01\% & +1.9\% & +1.1\%\\
& Gaussian & 1.01 / 1.06 & 0.02\% & +1.9\% & +1.1\%\\
\bottomrule
\end{tabular}
\end{table}

We use the calibration rule from \cite{wu2026selfsizing} as the baseline. For $k=3$, that work groups results by first-round size and takes the worst 1\% quantile $q_{0.01}$ of post-failure $\widehat d/d$ over several loads and sign compositions (0.8045, 0.8997, 0.9482 for $M_1=256,1024,4096$). We use its upper-bound component $U=\lceil\widehat d/q_{0.01}\rceil$ with the same second-round table length as our two rules. The comparison excludes that protocol's capacity coefficient, deduction of elements recovered in the first round, and joint peeling across the two rounds. Table~\ref{tab:vldb} compares this rule under both encodings, with and without accounting for $|b|$.

\begin{table}[!htbp]
\centering
\caption{Calibration rule versus the Skellam rule ($k=3$, $u=0.01$). The conditional miss rate /u gives the median / maximum for balanced configurations and the range for one-sided configurations. Final failure is the maximum. Total bytes are the median relative to the oracle.}
\label{tab:vldb}
\small
\begin{tabular}{@{}rlcccc@{}}
\toprule
$M_1$ & rule & balanced & one-sided & final failure & total bytes\\
\midrule
256 & standard encoding, calibration factor & 0.92 / 1.02 & 0.09--0.90 & 0.27\% & +13.9\%\\
& standard encoding, calibration factor, with $|b|$ & 0.92 / 1.02 & 0 & 0.26\% & +13.9\%\\
& random sign, calibration factor & 0.87 / 0.97 & 0.85--0.93 & 0.23\% & +13.9\%\\
& random sign, Skellam rule & 0.72 / 0.80 & 0 & 0.24\% & +14.6\%\\
1024 & standard encoding, calibration factor & 0.88 / 0.94 & 0.05--1.03 & 0.07\% & +6.4\%\\
& standard encoding, calibration factor, with $|b|$ & 0.88 / 0.94 & 0 & 0.07\% & +6.4\%\\
& random sign, calibration factor & 0.98 / 1.03 & 0.80--1.04 & 0.08\% & +6.4\%\\
& random sign, Skellam rule & 0.89 / 0.92 & 0 & 0.06\% & +6.5\%\\
4096 & standard encoding, calibration factor & 0.84 / 0.86 & 0.36--0.89 & 0.02\% & +3.1\%\\
& standard encoding, calibration factor, with $|b|$ & 0.84 / 0.86 & 0 & 0.02\% & +3.1\%\\
& random sign, calibration factor & 0.90 / 0.94 & 0.77--0.95 & 0.02\% & +3.1\%\\
& random sign, Skellam rule & 0.97 / 1.01 & 0 & 0.01\% & +3.1\%\\
\bottomrule
\end{tabular}
\end{table}

With the same calibration factor, the conditional miss rate on one-sided configurations ranges from $0.05u$ to $1.03u$ under standard encoding and from $0.77u$ to $1.04u$ under random-sign encoding. At fixed load 0.95--1.0, the rates for one-sided and balanced differences under standard encoding are $0.05$--$0.40u$ and $0.86$--$0.94u$, respectively. Under random-sign encoding, they are $0.80$--$0.95u$ and $0.85$--$1.01u$. Accounting for $|b|$ gives zero one-sided miss under both encodings. In paired comparisons by trial, with the calibration factor and the choice of whether to account for $|b|$ fixed, the success-rate difference between random-sign encoding and standard encoding lies within $\pm1\times10^{-3}$. Intervals exclude zero for 2 of the 30 comparisons, both at $M_1=256$ and load 1.0. The median total-byte difference is about 0.01\%.

\subsection{Computation Cost}\label{app:timing}
For the computation-cost measurement (\exptag{COST}), we use a single-threaded Go implementation on a dual-socket Linux host with Intel Xeon Silver 4410Y processors (nominal 2.0 GHz, 24 physical cores and 48 logical CPUs). Before timing, we pause the host's task queue and wait for all tasks in the current batch to finish, then pin the program to one core. We record that core's frequency and system load every 10 seconds. The frequency remains at 2.001 GHz throughout the run.

We store Poisson tables in a main-memory cache managed by the program. Each cache entry corresponds to one value in a candidate $n$'s Poisson table and stores its probability and cumulative probability as two double-precision numbers, totaling 16 bytes. The cache holds 16,000,000 entries (about 256 MB). When the cache exceeds its limit, the program evicts tables in least-recently-used order until the occupancy plus the size of the incoming table is at most half the capacity, then inserts the new table. We time sizing under three cache conditions. In cold, each sizing starts with an empty cache. In session, the same receiver performs consecutive synchronizations at fixed $(M,k)$, each possibly involving one or two protocol rounds and each processing a new observation. The cache persists between synchronizations. In warm, we immediately recompute the same observation once.

For sizing timings, we use the first 200 failure records per configuration from the statistical experiments at $M=16384$. We repeat each of the four modes, cold, session, warm, and Gaussian rule, 10 times, rotating their order across repetitions. For protocol timings, we use $M_1=16384$ and 100,000 common elements. We repeat the same 40 trials per configuration 10 times and rotate the branch order across repetitions. We count first-round decoding time per trial. Figure~\ref{fig:sizing} and Table~\ref{tab:proto} report results from two runs on the same host with the same settings. The protocol results differ by less than 1\% from the first run.

Appendix~\ref{app:search} predicts sizing cost roughly proportional to $\lambda_d$, or approximately $c^2$. In Figure~\ref{fig:sizing}, $M$ is fixed and load is proportional to $c$. The cold time grows roughly as load to the power 1.9, or $c^{1.9}$, from 0.06--0.12 ms at load about 0.08 ($c=8$) to 69--84 ms at load 3. The session time is about one fifth of cold, and warm is about one sixth to one seventh. At load 3, the Poisson tables requested cumulatively in session exceed the cache limit. The program evicts some old tables, and the session time rises to 41--55 ms. Gaussian rule sizing takes 0.06--0.29 $\mu$s. In the full protocol, the total times for the Gaussian rule and oracle differ by less than 1\%. Peak session memory for Poisson tables is 49--86 MB at load 1.05 and 128--174 MB at load 1.5, reaching the 256 MB limit at load 3.

\begin{figure}[!htbp]
\centering\includegraphics[width=\linewidth]{../figure/fig-sizing.pdf}
\caption{Skellam rule time per sizing ($M=16384$, single thread, cache limit 16,000,000 entries). (a) cold; (b) session (solid) and warm (dashed).}
\label{fig:sizing}
\end{figure}

The decoder treats cells with count $\pm1$ as candidate pure cells and verifies them with the checksum. Under standard encoding with one-sided differences, all keys have the same sign. A cell with count $\pm1$ therefore contains exactly one key, and every candidate passes. With both positive and negative difference elements, three keys with signs $+1,+1,-1$, for example, can occupy the same cell and also give count $+1$. The checksum rejects such candidates. We measure pure-cell checks, checksum rejections, and sign hashes per recovered key in a separate run (30 trials per configuration, 3 repetitions). Random-sign encoding checks an average of 1.77 candidates ($k=3$) and 2.53 candidates ($k=4$). Standard encoding checks 1.00 under one-sided differences and matches random-sign encoding under balanced differences. The first-round decoding time ratio (random sign / standard) is 0.99--1.06 for balanced differences and 1.26--2.42 for one-sided differences at load 1.0--1.5. At load 3, the one-sided ratio is 3.1 because standard encoding stops almost immediately under one-sided differences. Both absolute times are only 0.04--0.16 ms, less than 1\% of total protocol time.

For large parameters, we can replace the exact Skellam CDF with the Gaussian approximation. Across the 33 configurations in the exploration phase with Poisson mean above 200, this changes $U$ for at most 1.8\% of samples and leaves the miss count unchanged. We use the exact rule for all reported results.
